%---------------------------------------------------------------------------------
%                西南交通大学研究生学位论文：致谢
%---------------------------------------------------------------------------------
\chapter*{致\qquad{}谢}
\addcontentsline{toc}{chapter}{致\qquad{}谢}
“竢实扬华，自强不息。” 岁月不居，时节如流。
值此母校西南交通大学建校一百三十周年之际，回首在交大度过的七载春秋，心中感慨万千。
这七年，不仅是我探求学术真知的求学历程，更是我人格重塑与理想扎根的关键阶段。
生逢双甲又十的荣光，何其有幸，能将青春最赤诚的岁月镌刻在天佑鸿哲之间。
\par
首先，我要向恩师郑宇教授致以最崇高的敬意与谢忱。
郑老师治学严谨、高屋建瓴，其敏锐的学术洞察力与前瞻性的技术视野，引领我步入科研的殿堂。
不仅如此，郑老师在为人处世上的言传身教，更为我构建了深刻的人生方法论。
感念郑老师在我迷茫之时的悉心指引与包容鼓励，这份润物无声的师恩，我将永远铭记，并在未来的职业生涯中以此自勉，砥砺前行。
\par
感谢何华均师兄和麻志鹏师兄。两位师兄既有深耕科研的定力，亦有青云直上的气概。
你们严谨克难的科研精神和精益求精的实践态度，潜移默化地影响着我。
感谢你们在我遇到瓶颈时提供的无私指导，让我求学路走得愈发踏实笃定。
\par
研途漫漫，感谢在京东科技度过的充实时光。 这里承载了我研究生生涯的大半记忆。
特别感谢韩博洋、孟垂实两位老师的悉心指导与包容，让我在工业界实战中飞速成长。
感谢张晓龙、王璐璐、赵大燕、黎世骄、陈海乐等师兄师姐的并肩作战与温情陪伴，你们的热忱让我的实习生活熠熠生辉。
感谢郭成杰、孙冉润、阙柱沪等同门以及吴晨宇、虞杰杰、巴聪聪等师弟在工作与科研中的默契协作。
\par
感恩在清华大学智能产业研究院以及字节跳动的宝贵经历。让我深刻体会到了“敢为极致、勇攀高峰”的精神。
这些经历不仅磨炼了我的韧性，更拓宽了我看世界的视野。
\par
特别致谢陪伴我走过七轮寒暑的同门挚友：徐梓航、崔智源、柴江波、王一凡。
我们的情谊起笔于那个疫前未染尘埃的仲秋，在拔尖班的初见中定格；
又在求职与毕业尘埃落定后的盛夏，于奔赴北上深杭的万水千山间，写下这一阶段的终章。
科学研究表明，人体细胞每七年就会完成一次整体的更新；
何其有幸，在这场漫长的更迭中，是你们陪我见证了从旧我向新我的重塑与蜕变。
其深厚羁绊，不仅是学术路上的并肩，更是我求学岁月里最温情、也最坚韧的收获。
\par
最后，感谢一路上相知相遇的每一个人，恕于篇幅不一一列举。

感谢父母家人二十多年以来的坚定支持与相信。感谢那个在多地奔波却从未言弃、在深夜坚守却依然满怀勇气的自己。

一百三十载峥嵘岁月，母校常新；七年砥砺前行，初心不改。

